
\documentclass{article}
\usepackage{colortbl}
\usepackage{makecell}
\usepackage{multirow}
\usepackage{supertabular}

\begin{document}

\newcounter{utterance}

\centering \large Interaction Transcript for game `cladder', experiment `full\_v1.5\_default', episode 10063 with Dialogue Architects\_\_qwen3.5{-}9b\_\_SCoRe\_Qwen3.5{-}9B.
\vspace{24pt}

{ \footnotesize  \setcounter{utterance}{1}
\setlength{\tabcolsep}{0pt}
\begin{supertabular}{c@{$\;$}|p{.15\linewidth}@{}p{.15\linewidth}p{.15\linewidth}p{.15\linewidth}p{.15\linewidth}p{.15\linewidth}}
    \# & \multicolumn{2}{c}{Player} && \multicolumn{2}{c}{Game Master} \\
    \hline

    \theutterance \stepcounter{utterance}  
    & & & \multicolumn{4}{p{0.6\linewidth}}{
        \cellcolor[rgb]{0.9,0.9,0.9}{
            \makecell[{{p{\linewidth}}}]{
                \texttt{\tiny{[P1$\langle$GM]}}
                \texttt{You are an expert in causal inference. The following question is not a typical commonsense query, but rather a meticulously designed question created by a professor specializing in causal inference, intended to assess the students' mastery of the course content.} \\
\texttt{Imagine a self{-}contained, hypothetical world with only the following conditions, and without any unmentioned factors or causal relationships: Zory has a direct effect on zuph and glimx. Zuph has a direct effect on jyka. Jyka has a direct effect on glimx. Zory is unobserved. We know that zory causes zuph. zuph causes not jyka. zory or jyka causes glimx. We observed an individual is zory. Would an individual is not glimx if not zuph instead of zuph?} \\
\texttt{Start your answer with 'Yes' or 'No', followed by additional reasoning or evidence to support your explanation.} \\
            }
        }
    }
     \\ \\

    \theutterance \stepcounter{utterance}  
    & \multicolumn{4}{p{0.6\linewidth}}{
        \cellcolor[rgb]{0.561,0.737,0.561}{
            \makecell[{{p{\linewidth}}}]{
                \texttt{\tiny{[P1$\rangle$GM]}}
                \texttt{No} \\
\\ 
\texttt{**Reasoning:**} \\
\\ 
\texttt{To determine the answer, we must construct the causal graph based strictly on the provided conditions and evaluate the counterfactual query: "Would an individual be not glimx if not zuph instead of zuph?" given that the individual is observed to be zory.} \\
\\ 
\texttt{**1. Constructing the Causal Graph:**} \\
\texttt{Based on the text, the direct causal relationships are:} \\
\texttt{*   \$Zory \textbackslash{}rightarrow Zuph\$ (Explicitly stated: "zory causes zuph")} \\
\texttt{*   \$Zuph \textbackslash{}rightarrow Jyka\$ (Explicitly stated: "Zuph has a direct effect on jyka")} \\
\texttt{*   \$Jyka \textbackslash{}rightarrow Glimx\$ (Explicitly stated: "Jyka has a direct effect on glimx")} \\
\texttt{*   \$Zory \textbackslash{}rightarrow Glimx\$ (Explicitly stated: "Zory has a direct effect on... glimx")} \\
\\ 
\texttt{*Note on Contradictions:* The prompt contains a contradiction regarding the effect of Zuph on Jyka ("Zuph has a direct effect on jyka" vs. "zuph causes not jyka"). In causal inference problems of this type, the structural definition ("has a direct effect") usually establishes the existence of the arrow, while the specific outcome description ("causes not") describes the functional form or a specific realization. However, the most critical structural element for this specific counterfactual is the path from Zory to Glimx. The prompt explicitly states: "Zory has a direct effect on... glimx." This creates a direct causal link \$Zory \textbackslash{}rightarrow Glimx\$ that exists independently of the \$Zory \textbackslash{}rightarrow Zuph \textbackslash{}rightarrow Jyka \textbackslash{}rightarrow Glimx\$ path.} \\
\\ 
\texttt{**2. Analyzing the Counterfactual:**} \\
\texttt{*   **Factual World (\$M\$):** We observe \$Zory = 1\$ (true).} \\
\texttt{    *   Since \$Zory \textbackslash{}rightarrow Zuph\$, then \$Zuph = 1\$.} \\
\texttt{    *   Since \$Zuph = 1\$, \$Jyka\$ is affected (let's assume \$Jyka = 1\$ for the sake of the path, though the exact value depends on the functional form).} \\
\texttt{    *   Since \$Zory = 1\$ and there is a direct effect \$Zory \textbackslash{}rightarrow Glimx\$, \$Glimx\$ is influenced by Zory.} \\
\texttt{    *   The query asks about the state of \$Glimx\$ in the counterfactual world.} \\
\\ 
\texttt{*   **Counterfactual World (\$M\_\{Zuph=0\}\$):** We intervene to set \$Zuph = 0\$ while keeping the observed fact \$Zory = 1\$.} \\
\texttt{    *   **Path 1 (Blocked):** The path \$Zory \textbackslash{}rightarrow Zuph \textbackslash{}rightarrow Jyka \textbackslash{}rightarrow Glimx\$ is broken because \$Zuph\$ is set to 0. Consequently, the influence of Zory on Jyka (and subsequently Glimx via Jyka) is severed.} \\
\texttt{    *   **Path 2 (Active):** The prompt explicitly states: "Zory has a direct effect on... glimx." This implies a direct causal arrow \$Zory \textbackslash{}rightarrow Glimx\$.} \\
\texttt{    *   In the counterfactual world, \$Zory\$ is still observed to be true (\$Zory = 1\$). Because the direct effect \$Zory \textbackslash{}rightarrow Glimx\$ exists and is not mediated by Zuph, the value of \$Zory\$ will still cause a change in \$Glimx\$.} \\
\\ 
\texttt{**3. Evaluating the Outcome:**} \\
\texttt{The question asks: "Would an individual be **not glimx** if not zuph...?"} \\
\texttt{This is asking if \$Glimx = 0\$ (not glimx) is the result of the intervention.} \\
\\ 
\texttt{*   In the factual world, \$Glimx\$ is likely 1 (assuming the direct effect of Zory is positive/activating, which is the standard interpretation of "has a direct effect" in binary contexts unless specified otherwise).} \\
\texttt{*   In the counterfactual world (\$Zuph=0\$), the path through Jyka is cut. However, the direct path \$Zory \textbackslash{}rightarrow Glimx\$ remains active because \$Zory\$ is still 1.} \\
\texttt{*   Unless the prompt implies that the *only* way Zory affects Glimx is *through* Zuph (which contradicts the explicit statement "Zory has a direct effect on... glimx"), the intervention on Zuph does not eliminate the effect of Zory on Glimx.} \\
\texttt{*   Therefore, \$Glimx\$ will likely still be 1 (glimx), not 0 (not glimx).} \\
\\ 
\texttt{Even if we consider the contradictory phrase "zuph causes not jyka" as defining the functional form where \$Zuph=1 \textbackslash{}implies Jyka=0\$, the logic holds:} \\
\texttt{1.  Factual: \$Zory=1 \textbackslash{}implies Zuph=1 \textbackslash{}implies Jyka=0\$. But \$Zory=1 \textbackslash{}implies Glimx=1\$ (direct effect). So \$Glimx=1\$.} \\
\texttt{2.  Counterfactual: \$Zory=1\$, but we force \$Zuph=0\$.} \\
\texttt{    *   Does \$Zuph=0\$ change \$Jyka\$? If \$Zuph\$ was the *only* cause of Jyka, then \$Jyka\$ might become 1 (or 0 depending on default).} \\
\texttt{    *   Does this change \$Glimx\$? \$Glimx\$ is caused by \$Zory\$ (direct) and \$Jyka\$.} \\
\texttt{    *   Since \$Zory=1\$ directly causes \$Glimx\$, \$Glimx\$ remains 1 regardless of the state of \$Jyka\$ (unless there is a specific inhibitory logic not stated, e.g., "Glimx is true ONLY if Jyka is true," which is not stated). The prompt says "Zory or Jyka causes glimx," implying an inclusive OR or additive effect.} \\
\\ 
\texttt{Since \$Zory\$ directly causes \$Glimx\$, setting \$Zuph\$ to false does not negate the effect of \$Zory\$ on \$Glimx\$. Therefore, the individual would still be **glimx**, not **not glimx**.} \\
\\ 
\texttt{The statement "Would an individual be not glimx" is false.} \\
\\ 
\texttt{**Conclusion:**} \\
\texttt{No, the individual would not be "not glimx" (i.e., they would still be glimx) because Zory has a direct effect on Glimx independent of Zuph. Intervening to set Zuph to false only blocks the indirect path through Jyka, but the direct path from Zory to Glimx remains active.} \\
            }
        }
    }
    & & \\ \\

    \theutterance \stepcounter{utterance}  
    & & & \multicolumn{2}{p{0.3\linewidth}}{
        \cellcolor[rgb]{0.9,0.9,0.9}{
            \makecell[{{p{\linewidth}}}]{
                \texttt{\tiny{[GM$|$GM]}}
                \texttt{No} \\
            }
        }
    }
    & & \\ \\

    \theutterance \stepcounter{utterance}  
    & & & \multicolumn{2}{p{0.3\linewidth}}{
        \cellcolor[rgb]{0.9,0.9,0.9}{
            \makecell[{{p{\linewidth}}}]{
                \texttt{\tiny{[GM$|$GM]}}
                \texttt{no} \\
            }
        }
    }
    & & \\ \\

    \theutterance \stepcounter{utterance}  
    & & & \multicolumn{2}{p{0.3\linewidth}}{
        \cellcolor[rgb]{0.9,0.9,0.9}{
            \makecell[{{p{\linewidth}}}]{
                \texttt{\tiny{[GM$|$GM]}}
                \texttt{game\_result = WIN} \\
            }
        }
    }
    & & \\ \\

\end{supertabular}
}

\end{document}
